\documentclass[runningheads]{llncs}

\usepackage[utf8]{inputenc}
\usepackage[T1]{fontenc}
\usepackage[english]{babel}
\usepackage{booktabs}
\usepackage{hyperref}
\usepackage{amsmath}
\usepackage{array}
\usepackage{float}
\usepackage{graphicx}

\hypersetup{colorlinks=true, linkcolor=black, citecolor=black, urlcolor=blue}

\begin{document}

\title{Multi-Label ICD-10 Classification of Greek Cardiology Discharge Summaries}
\subtitle{System Description for the ElCardioCC 2026 Shared Task}

\author{
    Nikolaos Strafiotis \and
    Panteleimon Stanimeros \and
    Vasiliki Katsara \and
    Georgios Chalkias \and
    Glykeria Tsavlidou \and
    Efrosyni Nalmpanti \and
    Panteleimon Stamatakis \and
    Ioannis Refanidis\thanks{Corresponding author.}
}

\authorrunning{Strafiotis et al.}

\institute{
    Department of Applied Informatics, University of Macedonia, Thessaloniki, Greece\\
    MSc in Artificial Intelligence \& Data Analytics\\
    \email{\{aid26012, aid26006, aid26009, aid26004, aid26005, aid26014, aid26013\}@uom.edu.gr}\\
    \email{yrefanid@uom.edu.gr}
}

\maketitle

% ─────────────────────────────────────────────────────────────────────────────
\begin{abstract}
We present our system for the ElCardioCC 2026 shared task on automatic ICD-10 coding
of Greek cardiology discharge summaries.
The system combines six heterogeneous components --- a fine-tuned Greek BERT model,
two XLM-RoBERTa variants (base and large), a hybrid information retrieval module, a
dictionary-based rule system, and a named entity recognition with entity linking
pipeline --- under a declarative metaheuristic ensemble framework.
Individual models are trained with task-specific loss functions (Asymmetric Loss,
ZLPR) and per-label threshold tuning on the validation set.
The ensemble fuses outputs through ten base strategies and four composition operators,
with weights searched via random restarts, hill-climbing, and Variable Neighbourhood
Search (VNS).
Our best submission achieves a micro-F1 of \textbf{0.8667} on the official test set,
improving by approximately 1.8 F1 points over the best standalone model.

\keywords{ICD-10 coding \and clinical NLP \and multi-label classification \and
BERT \and ensemble methods \and Greek language}
\end{abstract}

% ─────────────────────────────────────────────────────────────────────────────
\section{Introduction}
\label{sec:intro}

Clinical coding assigns standardised diagnostic codes to free-text medical records,
a mandatory step in hospital administration, epidemiological surveillance, and
insurance reimbursement. Manual coding is slow, costly, and prone to inter-annotator
disagreement, motivating the development of automatic approaches \cite{dong2022automated}.

The \textbf{ElCardioCC 2026} shared task targets this problem for Greek cardiology
discharge summaries. Systems must predict all applicable ICD-10 codes from the WHO
10th revision for each document --- a multi-label classification problem over 115 codes.

Three key challenges define the setting. First, \textit{language}: Greek is a
morphologically rich, low-resource language for clinical NLP. Discharge summaries
blend formal Greek medical terminology with transliterated English abbreviations such
as \textit{PCI}, \textit{TAVI}, \textit{CABG}, and \textit{ACS}, creating a
heterogeneous vocabulary that strains standard tokenisers. Second, \textit{document
length}: some summaries exceed the 512-token context limit of standard transformer
encoders, requiring explicit long-document strategies. Third, \textit{label imbalance}:
the 115 ICD-10 codes span frequencies from a single training instance to over 1,000,
with the majority of the label space being rare. A system that ignores rare labels
may still report a high macro-averaged F1 while producing clinically useless
predictions for uncommon diagnoses.

Our approach is driven by the observation that no single model architecture dominates
across all labels and all document types. Frequent labels are best handled by
distributional neural models, while procedural codes with predictable lexical surface
forms are more precisely covered by dictionary-based rules. Rare labels benefit from
high-recall retrieval-based evidence, and mention-level context is critical for
disambiguation between closely related codes (e.g.\ I21 vs.\ I22 vs.\ I25). We
therefore build a system that combines all four paradigms and uses a metaheuristic
ensemble to find the best label-level fusion strategy on the validation set.

The remainder of this paper is structured as follows. Section~\ref{sec:background}
introduces key background concepts. Section~\ref{sec:task} describes the shared task
and dataset. Section~\ref{sec:method} presents the full system methodology.
Section~\ref{sec:results} reports experimental results. Section~\ref{sec:conclusion}
concludes.

% ─────────────────────────────────────────────────────────────────────────────
\section{Background}
\label{sec:background}

\subsection{Transformer Language Models for Clinical NLP}

The Transformer architecture \cite{vaswani2017attention} underlies most state-of-the-art
NLP systems. It processes text through stacked self-attention layers, learning to
attend selectively to relevant context across the full input sequence. BERT
\cite{devlin2019bert} pre-trains a deep bidirectional Transformer using masked language
modelling on large corpora, producing contextual representations that transfer
effectively to downstream tasks with limited labelled data.

Clinical NLP has widely adopted BERT-style models. Domain-specific variants pre-trained
on medical text (BioBERT, ClinicalBERT) consistently outperform general-purpose models
on medical tasks \cite{lee2020biobert}. For Greek clinical text, we use
\texttt{nlpaueb/bert-base-greek-uncased-v1} (Greek BERT) \cite{greekbert}, pre-trained
on Wikipedia, European Parliament proceedings, and web text with a 35K Greek vocabulary.

XLM-RoBERTa \cite{conneau2020unsupervised} extends the RoBERTa training recipe
\cite{liu2019roberta} (longer training, larger batches, no next-sentence prediction)
to 100 languages using a 250K multilingual SentencePiece vocabulary. Its multilingual
vocabulary naturally handles the Greek--Latin code-switching present in clinical
abbreviations, and the large variant (560M parameters) provides greater representational
capacity than BERT-base.

\subsection{Multi-Label Classification and Loss Functions}

In multi-label classification each instance is assigned any subset of a fixed label
set. The standard approach adds a sigmoid-activated linear head over the transformer
output, producing per-label probabilities $p_i \in [0,1]$, and applies a threshold
$t_i$ per label: predict label $i$ if $p_i > t_i$.

Label imbalance is the central training challenge. For a dataset with $L$ labels,
each training example is positive for a small subset and negative for all others.
Standard binary cross-entropy is dominated by easy negatives. Two loss functions
address this.

\textbf{Asymmetric Loss (ASL)} \cite{ridnik2021asymmetric} applies asymmetric focal
down-weighting with a hard factor for easy negatives and a soft factor for positives:
\begin{equation}
\mathcal{L}_{\text{ASL}} = \begin{cases}
(1-p)^{\gamma^+}\log p & y=1 \\
p_m^{\gamma^-}\log(1-p_m) & y=0
\end{cases}
\end{equation}
where $p_m = \max(p-m,0)$ shifts the probability of negatives downward by margin $m$,
with $\gamma^-=4$, $\gamma^+=1$, and $m=0.05$ in our configuration.

\textbf{ZLPR} (Zero-Logit Paired Ranking) loss is a focal-style ranking objective that
directly optimises the margin between positive and negative label logits, without
requiring calibrated probability estimates. This makes it effective when the neural
model is used mainly for ranking rather than probability output.

\subsection{Information Retrieval}

\textbf{BM25} (Okapi BM25) is a classical probabilistic sparse retrieval function that
scores candidates by term overlap with the query using inverse document frequency
weighting and term-frequency saturation.

\textbf{Dense retrieval} \cite{reimers2020making} encodes queries and candidates into
a shared embedding space using a neural encoder, then retrieves by cosine similarity.
Multilingual sentence transformers such as \texttt{paraphrase-multilingual-MiniLM-L12-v2}
are effective across 50+ languages without language-specific fine-tuning.

\textbf{Reciprocal Rank Fusion (RRF)} \cite{cormack2009reciprocal} merges multiple
ranked lists without requiring score normalisation:
$\text{RRF}(d) = \sum_i w_i / (k + \text{rank}_i(d))$ with constant $k=60$.

\subsection{Metaheuristic Optimisation}

Metaheuristics search large solution spaces without gradient information.

\textbf{Hill-climbing with random restarts} iteratively applies small perturbations,
accepting improvements greedily. Multiple random starting points reduce the risk of
getting stuck in poor local optima.

\textbf{Variable Neighbourhood Search (VNS)} \cite{mladenovic1997variable} extends
hill-climbing by systematically varying the neighbourhood structure: when no improvement
is found in the current neighbourhood, a ``shake'' perturbation of increasing size is
applied before restarting local search. This allows VNS to escape local optima without
the full cost of random restarts, and is particularly well suited to non-smooth
objectives such as the discrete micro-F1 metric.

% ─────────────────────────────────────────────────────────────────────────────
\section{Task and Dataset}
\label{sec:task}

\subsection{Task Definition}

ElCardioCC 2026 is a multi-label ICD-10 coding shared task for Greek cardiology
discharge summaries, organised as part of the BioASQ/CLEF 2026 evaluation campaign.
The label set contains 115 ICD-10 codes drawn from cardiology practice in Greece.

The primary evaluation metric is a \textbf{relaxed micro-F1}: semantically equivalent
codes are grouped into clinical clusters, so a prediction from the same cluster as the
gold label is not fully penalised. This reflects clinical reality, where closely related
codes (e.g.\ different subtypes of atrial fibrillation) carry similar treatment
implications. Micro-averaging ensures that frequent codes dominate the score, consistent
with their greater clinical impact.

\subsection{Dataset}

The dataset consists of 2,500 anonymised discharge summaries split into training (2,000),
validation (250), and test (250) sets, with a separate blind test set held by the
organisers. Documents are moderately long: mean 280 tokens ($\approx$2,000 characters),
maximum 1,218 tokens. Only 2.4\% of documents exceed the 512-token BERT limit.

\textbf{Label distribution.} The 115 codes are highly skewed. The five most frequent
(Z95, I25, I21, I48, E11) account for $\sim$41\% of all annotations. At the other
extreme, 30 codes appear fewer than 10 times and 4 codes appear only once in the entire
dataset (Table~\ref{tab:label_stats}). This long-tailed distribution severely challenges
standard classifiers and requires explicit strategies for rare-label coverage.

\begin{table}[h]
\centering
\caption{Label frequency distribution across 115 ICD-10 codes.}
\label{tab:label_stats}
\begin{tabular}{lrr}
\toprule
Frequency band & \# codes & Approx.\ \% of instances \\
\midrule
$\geq 500$  & 5  & 41\% \\
100--499    & 22 & 36\% \\
10--99      & 58 & 21\% \\
$< 10$      & 30 & $<$2\% \\
\bottomrule
\end{tabular}
\end{table}

\subsection{Preprocessing and Data Augmentation}

A shared preprocessing pipeline is applied across all splits: (1) Unicode
normalisation (NFD) and accent stripping; (2) lowercasing for uncased models
(Greek BERT only); (3) character filtering to Greek/Latin/digits/punctuation;
(4) whitespace collapse. For XLM-RoBERTa (case-sensitive SentencePiece tokeniser),
only whitespace normalisation is applied to preserve abbreviation signals such as
``PCI'' and ``TAVI''.

Multi-label stratified splitting with patient-level grouping ensures that all records
from the same patient belong to the same split, preventing data leakage. Rare labels
(fewer than 2 instances) are assigned via forced placement to guarantee at least one
training example per class.

Synonym-based data augmentation is applied to the XLM-RoBERTa base training track.
Dictionary terms are swapped with synonyms at probability 0.35. To avoid amplifying
noise from very rare labels, augmentation multipliers are frequency-adaptive:
$\times$4 for labels in the 30--99 instance range, $\times$3 for 100--199, $\times$2
for 200--299, and no augmentation outside this range.

% ─────────────────────────────────────────────────────────────────────────────
\section{Methodology}
\label{sec:method}

The system pipeline is: (i) train six independent components; (ii) tune per-model
decision thresholds on the validation set; (iii) search for the best ensemble fusion
strategy. Figure~\ref{fig:pipeline} provides an overview.

\begin{figure}[h]
\centering
\small
\renewcommand{\arraystretch}{1.2}
\begin{tabular}{c}
\textbf{Raw JSONL} (train / validation / test) \\
$\downarrow$ \\
\textbf{Preprocessing \& Stratified Split} \\
$\downarrow$ \\
\begin{tabular}{|c|c|c|c|c|c|}
\hline
Greek & XLM-R & XLM-R & Hybrid & Dictionary & NER \\
BERT & Large & Base & IR & Baseline & +EL \\
\hline
\end{tabular}
\\
$\downarrow$ \\
\textbf{Per-model threshold tuning} (validation set) \\
$\downarrow$ \\
\textbf{Metaheuristic Ensemble} (10 strategies + 4 compositions) \\
$\downarrow$ \\
\textbf{Final Predictions}
\end{tabular}
\caption{High-level system pipeline.}
\label{fig:pipeline}
\end{figure}

\subsection{Greek BERT Multi-Label Classifier}
\label{subsec:bert}

\subsubsection{Architecture.}
We fine-tune Greek BERT with a multi-label classification head consisting of: mean
pooling over all token representations (masked by the attention mask), Dropout (0.3),
Linear $768 \to 384$ with GELU activation and Dropout (0.3), and Linear $384 \to 115$.
Mean pooling over all tokens outperformed CLS-token extraction by $\sim$0.8 F1 points
on our validation set, consistent with findings that diagnostic information is
distributed throughout a discharge summary rather than concentrated in the opening
sentence. The maximum input sequence length is 256 tokens.

\subsubsection{Training configuration.}
Table~\ref{tab:bert_params} lists all training hyperparameters.

\begin{table}[h]
\centering
\caption{Greek BERT training hyperparameters.}
\label{tab:bert_params}
\begin{tabular}{ll}
\toprule
Hyperparameter & Value \\
\midrule
Base model        & \texttt{nlpaueb/bert-base-greek-uncased-v1} \\
Max length        & 256 tokens \\
Epochs / patience & 30 / 7 \\
Effective batch   & 32 (batch 16, grad.\ accum.\ 2) \\
Learning rate     & $3.0 \times 10^{-5}$ \\
LLRD factor       & 0.9 per encoder layer \\
Warmup            & 6\% of total steps \\
Weight decay      & 0.02 \\
Max grad norm     & 1.0 \\
Loss              & ASL ($\gamma^-=4$, $\gamma^+=1$, clip $=0.05$) \\
Precision         & FP16 mixed precision \\
Seed              & 42 \\
\bottomrule
\end{tabular}
\end{table}

\textbf{Layer-wise learning rate decay (LLRD)} applies a multiplicative factor of 0.9
per encoder layer from the head downward. The effective LR for layer $l$ (counting
from the output) is $\text{LR} \times 0.9^l$, so lower layers that encode general
Greek language structure are updated more conservatively than the upper task-specific
layers. LLRD improved validation micro-F1 by $\approx$0.4 points, with the most gain
on rare labels where lower-layer representations carry more predictive weight.

\subsubsection{Threshold tuning.}
After training, per-label decision thresholds are tuned in two passes.
\textit{Pass 1:} a global sweep over $t \in [0.45, 0.75]$ (step 0.01) identifies the
best single threshold across all 115 labels.
\textit{Pass 2:} each label is swept independently in the same range; per-label values
are accepted only for labels with $\geq$15 positive validation examples and only if
the improvement exceeds 0.001 F1 (to avoid overfitting on sparse data).
The resulting per-label thresholds range from 0.54 to 0.75: frequent labels converge
to lower thresholds reflecting higher model confidence, while rare labels require
higher thresholds due to systematic under-confidence. This tuning adds $\approx$1.5
F1 points over a fixed threshold of 0.5. Best validation micro-F1: \textbf{0.8165}.

\subsubsection{Training variants.}
Multiple variants were explored across development phases: an early baseline with
longer input sequences (\texttt{p2\_long\_asl\_warmlr}), a mid-regularisation variant
with lower weight decay (\texttt{p4\_asl\_mid\_reg}), and the final lean-head winner
(\texttt{p4\_winner\_lean\_head}). A 100-epoch variant of the winner was also trained
without early stopping and submitted as a standalone system.

\subsection{XLM-RoBERTa Classifiers}
\label{subsec:xlmr}

XLM-RoBERTa's 250K multilingual vocabulary handles Greek--Latin code-switching better
than Greek BERT's 35K Greek-only vocabulary, and the large variant (24 layers,
1024 hidden) provides greater representational capacity.

\subsubsection{Long-document strategy.}
A sliding window with stride 256 processes documents exceeding 512 tokens. One chunk
is randomly selected per document during training; all chunk logits are fused at
inference:
\begin{equation}
\hat{y} = 0.5 \cdot \bar{y}_{\mathrm{mean}} + 0.5 \cdot y_{\mathrm{max}}
\end{equation}
The mean term captures the overall label distribution; the max term ensures strong
local signals from any single chunk are not averaged away.

\subsubsection{Large variant.}
Uses ZLPR loss, batch size 2, gradient accumulation $\times$16 (effective batch 32),
LR $2\times10^{-5}$, weight decay 0.05, classifier dropout 0.15, 40 epochs with early
stopping (patience 5). The small physical batch size was imposed by GPU memory
constraints given the 560M parameter model.

\subsubsection{Base variant.}
Uses ASL, batch size 16, and optional K-fold cross-validation for generating
fold-averaged predictions. It provides a computationally lighter alternative to the
large variant and contributes additional diversity to the ensemble.

\subsection{Information Retrieval Module}
\label{subsec:ir}

The IR module frames coding as a retrieval problem: rank ICD-10 codes by similarity to
the discharge summary. This requires no labelled training data, generalises to rare
codes, and provides high recall at the cost of precision.

Three retrieval channels are combined by RRF ($k=60$, equal weights $w_i=1$,
candidate pool $\max(6 \cdot \text{top\_k},\, 80)$):
\begin{itemize}
  \item \textbf{BM25:} sparse term-overlap ranking over ICD-10 code descriptions
  \item \textbf{TF-IDF:} cosine similarity in a TF-IDF vector space
  \item \textbf{Dense:} cosine similarity using \texttt{paraphrase-multilingual-MiniLM-L12-v2}
        embeddings (384 dimensions, L2-normalised)
\end{itemize}

Prediction: codes with RRF score $\geq 0.22 \times \text{top score}$, capped at 12
per document. The cap is critical: relaxing it from 12 to 25 drops F1 from 0.45 to
0.14 (Table~\ref{tab:ir_ablation}) by flooding precision with incorrect candidates.

\begin{table}[h]
\centering
\caption{IR module tuning results on the validation set.}
\label{tab:ir_ablation}
\begin{tabular}{llccc}
\toprule
Configuration & max\_codes & Recall & Precision & F1 \\
\midrule
Plain top-25  & 25 & 0.374 & 0.086 & 0.140 \\
Tuned hybrid  & 12 & 0.828 & 0.310 & \textbf{0.451} \\
\bottomrule
\end{tabular}
\end{table}

\subsection{Dictionary Baseline}
\label{subsec:dict}

A curated CSV (1,680 term-to-code entries) is matched by an Aho--Corasick automaton
with word-boundary enforcement and a 35-term false-positive blacklist. Cardiology
pattern rules add codes for named procedures and diagnoses (Table~\ref{tab:rules}).

\begin{table}[h]
\centering
\caption{Selected rule-based boosters (trigger terms translated from Greek).}
\label{tab:rules}
\begin{tabular}{ll}
\toprule
Trigger (Greek abbreviation / term) & Added codes \\
\midrule
PCI, PPCI, PTCA, angioplasty & Y84 \\
Stent, pacemaker & Z95 \\
TAVI & Y84, Z95, I35 \\
CABG & Y84, Z95, I25 \\
Myocarditis / pericarditis & I41 / I30 \\
Pulmonary embolism & I26 \\
Ventricular tachycardia & I47 \\
AV block, Mobitz II & I44 \\
ACS / acute coronary syndrome & I20, I21, I22, I24 \\
Acute renal failure / haemodialysis & N17 / Z99 \\
Respiratory failure & J96 \\
Amyloidosis & E85 \\
\bottomrule
\end{tabular}
\end{table}

Co-occurrence rules enforce clinical consistency: predicting I22 (subsequent MI) adds
I21 (acute MI); predicting I25 (chronic ischaemic disease) adds Z95 (vascular
implants). Optional filters for negation (4-token window) and family history sections
are available but disabled in the primary configuration.

Dictionary results on the training set: micro-F1 \textbf{0.7176} (precision 0.65,
recall 0.80), covering 99.76\% of documents and 73\% of label types. The 31 uncovered
labels include codes requiring inference from laboratory values or multi-sentence
implicit evidence.

\subsection{Named Entity Recognition and Entity Linking}
\label{subsec:nerel}

A Greek BERT BIO sequence tagger with a partial CRF layer identifies clinical mention
spans. Training: 6 epochs, batch 8, LR $2\times10^{-5}$, weight decay 0.01, early
stopping (patience 2), dynamic padding.

Each recognised mention is linked to ICD-10 codes via:
\begin{equation}
\text{score}(m,c) = 0.6 \cdot \text{prior}(m,c) + 0.4 \cdot \cos(\mathbf{e}_m,\mathbf{e}_c)
\end{equation}
where $\text{prior}(m,c)$ is a (mention, code) frequency counter from training
annotations, and $\mathbf{e}_m$, $\mathbf{e}_c$ are 200-character context and code
description embeddings from \texttt{paraphrase-multilingual-MiniLM-L12-v2}.
Best validation micro-F1: \textbf{0.7223} at epoch 5.

\subsection{Metaheuristic Ensemble}
\label{subsec:ensemble}

\subsubsection{Strategy space.}
The ensemble committee is the six components above. All strategies are defined in YAML,
enabling evaluation without code changes. \textbf{Ten base strategies} include:
weighted ensemble (learnable weights per component), per-label champion routing (select
highest-F1 model per label), per-label champion + vote, per-patient score routing
(select highest mean-confidence model per document), gated secondary (neural primary;
IR/NER fallback for uncertain documents), frequency-bucketed thresholding, dual-threshold
zones, top-K cutoff, and label correction (grid search on systematic (predicted, gold)
co-occurrence patterns).

\textbf{Four composed strategies} apply label-set algebra over base strategy outputs:
\textit{Union} (merge\_or: predict if any strategy predicts; maximises recall);
\textit{Intersection} (merge\_and: predict only if all agree; maximises precision);
\textit{$k$-of-$n$} (merge\_k2: predict if $\geq k$ strategies agree).

\subsubsection{Weight optimisation.}
Ensemble weights are optimised over the probability simplex ($\sum_i w_i = 1$,
$w_i \geq 0$) on the validation micro-F1 in two stages:

\textit{Stage 1 --- Random search + hill-climbing:} 10,000 evaluations with 2
random restarts. Candidate weight vectors are sampled uniformly from the simplex,
then locally refined by pairwise perturbations (increase one weight, decrease another)
accepted greedily if validation F1 improves.

\textit{Stage 2 --- VNS:} Starting from the hill-climbing solution, VNS iterates
shake--local-search cycles with increasing neighbourhood radius. When no improvement
is found in the current neighbourhood, the radius increases; when an improvement is
found, the radius resets. This escapes local optima that trap pure hill-climbing while
amortising the cost of random restarts.

The four submitted composed strategies are:
(1) per-label routing $\cup$ weighted;
(2) per-label champion+vote $\cup$ weighted;
(3) weighted $\cap$ correction (\textbf{best submission});
(4) 2-of-3 vote among weighted, per-label routing, and correction.

% ─────────────────────────────────────────────────────────────────────────────
\section{Experimental Results}
\label{sec:results}

\subsection{Individual Component Performance}

Table~\ref{tab:individual} reports standalone validation micro-F1 for each component.
Greek BERT is the strongest single model by a large margin. Dictionary and NER+EL
perform comparably; the IR module has the lowest F1 but the highest recall (0.83),
making it a key coverage contributor in the ensemble.

\begin{table}[h]
\centering
\caption{Individual component validation micro-F1.}
\label{tab:individual}
\begin{tabular}{lcc}
\toprule
Component & Val.\ micro-F1 & Notes \\
\midrule
Greek BERT (per-label thresholds) & 0.8165 & Best single model \\
NER + Entity Linking              & 0.7223 & Best at epoch 5 \\
Dictionary Baseline               & 0.7176 & 99.76\% document coverage \\
IR (hybrid RRF)                   & 0.4510 & Recall 0.83, Precision 0.31 \\
\bottomrule
\end{tabular}
\end{table}

\subsection{Ablation Studies}

\subsubsection{Threshold tuning.}
Table~\ref{tab:thr} shows the effect of progressively more granular thresholding.
Per-label tuning adds $\approx$1.5 F1 points over a fixed threshold of 0.5 and
$\approx$1.0 points over the global optimum.

\begin{table}[h]
\centering
\caption{Threshold tuning ablation for Greek BERT.}
\label{tab:thr}
\begin{tabular}{lc}
\toprule
Strategy & Val.\ micro-F1 \\
\midrule
Fixed $t=0.5$ & $\approx$0.79 \\
Global sweep $[0.45, 0.75]$ & $\approx$0.80 \\
Per-label sweep (2-pass) & \textbf{0.8165} \\
\bottomrule
\end{tabular}
\end{table}

\subsubsection{Impact of LLRD.}
Layer-wise learning rate decay improved validation F1 by $\approx$0.4 points. The
improvement is most pronounced for rare labels, where preserving pre-trained lower-layer
representations is more important than task-specific adaptation.

\subsubsection{Loss function comparison.}
ASL outperformed both binary cross-entropy and standard focal loss, primarily through
better recall on rare labels. BCE delivered good precision but lower recall; focal
loss improved recall relative to BCE but less than ASL due to its symmetric treatment
of positive and negative examples.

\subsubsection{Ensemble composition.}
Table~\ref{tab:comp} compares the three composition operators at a schematic level.
The intersection (\texttt{merge\_and}) achieves the best F1 by combining the weighted
ensemble's strong baseline with the correction module's precision improvement.

\begin{table}[h]
\centering
\caption{Ensemble composition operator comparison (validation set, schematic).}
\label{tab:comp}
\begin{tabular}{lccc}
\toprule
Strategy & Recall & Precision & F1 trend \\
\midrule
Union (merge\_or)          & highest & lowest  & moderate \\
Weighted ensemble          & mid     & mid     & strong \\
$k$-of-$n$ (merge\_k2)    & mid     & mid+    & strong \\
Intersection (merge\_and)  & mid$-$  & highest & \textbf{best} \\
\bottomrule
\end{tabular}
\end{table}

\subsection{Official Test Set Results}

Table~\ref{tab:results} reports all five submissions on the official test set. The best
submission uses the \texttt{merge\_and} composed strategy (weighted ensemble
intersected with the label correction module).

\begin{table}[h]
\centering
\caption{Official ElCardioCC 2026 test set results.}
\label{tab:results}
\begin{tabular}{lccc}
\toprule
Submission & Recall & Precision & F1 \\
\midrule
\textbf{ensemble (merge\_and: weighted + correction)} & \textbf{0.8510} & \textbf{0.8830} & \textbf{0.8667} \\
ensemble (merge\_k2: weighted + per-label + corr.)    & 0.8687 & 0.8539 & 0.8613 \\
ensemble (safer thr, merge\_k2)                       & 0.8767 & 0.8431 & 0.8596 \\
mlc\_greek\_bert\_100 (100-epoch standalone)          & 0.8194 & 0.8811 & 0.8491 \\
mlc\_greek\_bert (30-epoch standalone)                & 0.8310 & 0.8675 & 0.8489 \\
\bottomrule
\end{tabular}
\end{table}

The ensemble outperforms the best standalone model by \textbf{+1.8 F1 points}
(0.8667 vs.\ 0.8489), confirming the complementarity of the six components.
The 100-epoch Greek BERT achieves virtually the same F1 as the 30-epoch version
(0.8491 vs.\ 0.8489), indicating early convergence and the absence of significant
overfitting with extended training under ASL and weight decay.

Among ensemble submissions, the precision-oriented \texttt{merge\_and} yields the
highest overall F1 (0.8667), while \texttt{merge\_k2} variants achieve higher recall
at the cost of precision --- consistent with the 2-of-3 voting design. The safer
threshold variant of \texttt{merge\_k2} reaches the highest recall (0.8767) but
the lowest precision (0.8431).

\subsection{Error Analysis}

\textbf{Rare labels.} The 30 codes with $<$10 training instances are the primary error
source. For 10 of these, per-label thresholding is not applicable. The dictionary
rule-based system recovers some (e.g.\ E85 via the amyloidosis rule), but codes with
no lexical surface form (e.g.\ M35, R57) remain consistently unrecovered.

\textbf{Coverage gaps.} 31 label types lie outside the dictionary's coverage, including
metabolic (E00), haematological (D47), and musculoskeletal codes. These require
inference from laboratory values or implicit multi-sentence evidence.

\textbf{Label confusion.} I21/I22/I25 (acute MI, subsequent MI, chronic ischaemic
disease) frequently co-occur in a single document and require temporal reasoning to
assign correctly --- a capability beyond current bag-of-token models. Z95 (presence of
vascular implants) is the main false-positive source due to ubiquitous stent and
pacemaker mentions, motivating the \texttt{merge\_and} correction step.

\textbf{Calibration.} Greek BERT shows systematic under-confidence for rare labels:
true positive probabilities for rare codes often fall below 0.5. Per-label thresholding
mitigates this but becomes unreliable for labels with fewer than 15 positive validation
examples.

% ─────────────────────────────────────────────────────────────────────────────
\section{Conclusion}
\label{sec:conclusion}

We presented a multi-component system for automatic ICD-10 coding of Greek cardiology
discharge summaries, developed as part of the MSc in Artificial Intelligence and Data
Analytics at the University of Macedonia.

The system integrates Greek BERT and XLM-RoBERTa neural classifiers with a
dictionary-based rule system, a hybrid information retrieval module, and a NER+EL
pipeline. These are fused by a declarative metaheuristic ensemble that searches over
ten base strategies and four composition operators using random search, hill-climbing,
and Variable Neighbourhood Search. Our best submission achieves micro-F1
\textbf{0.8667}, a gain of 1.8 points over the standalone Greek BERT model.

The error analysis identifies three main limitations: rare-label under-coverage due to
insufficient training data and calibration issues; dictionary gaps for non-cardiology
comorbidity codes; and label confusion among temporally related ICD-10 codes.

Future work will address these through better calibration techniques for the long tail
(e.g.\ temperature scaling, isotonic regression), temporal reasoning modules for
distinguishing related codes (I21/I22/I25), larger Greek language models, and
cross-lingual transfer from well-annotated English clinical corpora.

% ─────────────────────────────────────────────────────────────────────────────
\section*{References}

\textit{[References will be added in a later version.]}

\smallskip\noindent The following works are cited in the text and will be formally
included:
\begin{itemize}
  \item Vaswani et al.\ (2017) --- Attention Is All You Need
  \item Devlin et al.\ (2019) --- BERT
  \item Liu et al.\ (2019) --- RoBERTa
  \item Conneau et al.\ (2020) --- XLM-RoBERTa
  \item Koutsikakis et al.\ (2020) --- Greek-BERT
  \item Ridnik et al.\ (2021) --- Asymmetric Loss for Multi-Label Classification
  \item Reimers \& Gurevych (2020) --- Making Monolingual Sentence Embeddings Multilingual
  \item Cormack et al.\ (2009) --- Reciprocal Rank Fusion
  \item Mladenovic \& Hansen (1997) --- Variable Neighbourhood Search
  \item Lee et al.\ (2020) --- BioBERT
  \item Dong et al.\ (2022) --- Automated clinical coding: a review
\end{itemize}

\end{document}
